\chapter*{Introducción}
\addcontentsline{toc}{chapter}{Introducción}
\markboth{Introducción}{Introducción}

El Método de los Elementos Finitos (MEF) es una de las herramientas centrales del análisis estructural y de la ingeniería computacional \autocite{zienkiewicz2013fem,cook2002concepts}, y su enseñanza forma parte del núcleo curricular de las carreras de ingeniería civil y mecánica \autocite[p.~1160]{perezsantiago2023fem}. Quien lo aprende debe articular tres planos que suelen presentarse por separado: la formulación matemática continua (elasticidad lineal, ecuaciones de equilibrio); su discretización mediante elementos isoparamétricos e integración numérica; y la interpretación física de los resultados, es decir, de los campos de desplazamiento, deformación y tensión (esfuerzo). Este trabajo desarrolla EduFEM, un software educativo de escritorio que reúne esos tres planos en un único entorno interactivo para problemas bidimensionales de elasticidad lineal estática, con elementos cuadriláteros de cuatro nodos (Q4) y de nueve nodos (Q9).

\section*{Planteamiento del problema}

La enseñanza del MEF tropieza con dos dificultades que se refuerzan entre sí. Por un lado, los textos clásicos del método, como los de Bathe y Reddy \autocite{bathe2014fem,reddy2006introduction}, desarrollan la teoría con un grado de abstracción considerable: las funciones de forma, la matriz Jacobiana, la matriz de deformación-desplazamiento y la integración de Gauss aparecen como pasos de una sucesión algebraica, y su conexión con la geometría concreta de un elemento no siempre es evidente para quien se inicia. Por otro lado, el software comercial de análisis opera como una caja negra: el usuario define geometría, material y cargas y obtiene resultados, pero el procesamiento intermedio que transforma esas entradas en tensiones permanece oculto \autocite[p.~1160]{perezsantiago2023fem}. El antecedente educativo más directo de este trabajo lo describía ya en 1998 en los mismos términos: con los programas comerciales, el estudiante queda limitado a introducir los datos del problema e interpretar los resultados, sin acceso a los detalles del análisis \autocite[p.~246]{suarez1998edelas2d}.

Esa opacidad tiene una consecuencia técnica concreta: un análisis cuyos resultados no pueden seguirse hasta sus datos ni comprobarse paso a paso. En la práctica profesional, un cálculo estructural se entrega con su memoria de cálculo precisamente para que un revisor pueda seguirlo y comprobarlo. El software comercial, en cambio, entrega los resultados del análisis por elementos finitos sin la memoria de su procedimiento interno: mirando el programa no es posible saber de dónde sale una tensión nodal, ni si las tensiones de los nodos cierran con las calculadas en el interior de los elementos. La pertinencia del problema para la formación está documentada. Un estudio de consenso entre expertos de la industria y de la academia reporta que la formación tiende a privilegiar las destrezas de operación sobre los conceptos fundamentales, y reclama que el egresado sepa planificar, verificar y validar sus análisis, no solo ejecutarlos \autocite[pp.~1159, 1171]{perezsantiago2023fem}.

Los recursos educativos existentes no cierran esa brecha para la elasticidad plana. Como muestra la \autoref{tab:comparativa}, ninguno de los recursos revisados reúne a la vez tres atributos: la exposición del procedimiento de cálculo completo sobre el modelo del usuario, una memoria de cálculo reproducible a mano y una verificación y validación publicada junto con la herramienta. A ello se suma que las herramientas comerciales no están diseñadas con un propósito didáctico explícito \autocite[p.~1160]{perezsantiago2023fem} y que su uso en el aula en idioma español es limitado. Esa es la \emph{situación problemática} de la que parte el trabajo: en la investigación tecnológica el problema no se elige libremente, sino que resulta de interpretar y diagnosticar una realidad concreta \autocite[p.~83]{garciacordoba2005tecnologica}.

En términos formales, el \textbf{problema científico} de esta investigación se formula como la siguiente interrogante: ¿de qué manera el uso del software de análisis estructural como caja negra podrá incidir en la baja trazabilidad y verificabilidad del análisis por el Método de los Elementos Finitos en elasticidad plana, en la Carrera de Ingeniería Civil de la Universidad Autónoma Tomás Frías, en la gestión 2026? Es la única formulación del problema que el documento emplea: la \autoref{sec:matriz-consistencia} la retoma textualmente y la \autoref{sec:validacion-hipotesis} la responde.

La variable causa es la forma en que el software expone el \emph{procedimiento de cálculo}, es decir, la secuencia de etapas que va de las funciones de forma a la recuperación de tensiones. Esa forma tiene dos niveles: como caja negra, que solo muestra datos y resultados, o con el procedimiento a la vista, etapa por etapa. La variable efecto es la trazabilidad y la verificabilidad del análisis. Se entiende por \emph{trazabilidad} la propiedad por la cual cada resultado puede seguirse hacia atrás, etapa por etapa, hasta los datos del modelo, con todas las magnitudes intermedias a la vista. Se entiende por \emph{verificabilidad} la propiedad por la cual cada magnitud, intermedia o final, puede comprobarse contra un patrón independiente —un cálculo manual, una solución analítica u otro programa— y esa comprobación da conforme dentro de una tolerancia fijada de antemano. Son las dos propiedades que la memoria de cálculo asegura en la práctica profesional. Ambas variables se observan en el software y en los casos de estudio, no en personas. La insuficiencia de una herramienta que exponga el procedimiento con esos atributos es la contradicción que origina el trabajo \autocite[p.~7]{alvarez_metodologia}.

El problema se delimita en cuatro planos. En lo disciplinar, a la elasticidad lineal plana con elementos cuadriláteros Q4 y Q9 (véase \emph{Alcance y limitaciones}). En lo institucional, a la formación de grado en ingeniería civil, con la Carrera de Ingeniería Civil de la Universidad Autónoma Tomás Frías como destinataria inmediata de la herramienta. En lo espacial, al aula y a la computadora personal del estudiante, sin requerir laboratorio de cómputo ni licencias. Y en lo temporal, al desarrollo, la verificación y la validación del software realizados durante la gestión 2026.

El \textbf{objeto de estudio} \autocite[p.~8]{alvarez_metodologia} es el análisis de medios continuos en elasticidad lineal plana por el Método de los Elementos Finitos: el procedimiento de cálculo que va del mapeo isoparamétrico a la recuperación de tensiones con elementos cuadriláteros Q4 y Q9, sobre el que se construye el motor, se miden las variables y se contrastan los resultados. El \textbf{campo de acción}, la parte de la situación problemática que el trabajo se propone mejorar \autocite[p.~13]{alvarez_metodologia}, son los medios de cálculo con que ese análisis se enseña en la Carrera: el software educativo que expone el procedimiento y la memoria de cálculo que lo documenta sobre el modelo del usuario.

\section*{Justificación}

El desarrollo de EduFEM se justifica en cuatro planos complementarios. En el plano académico, la herramienta acorta la distancia entre la formulación teórica del MEF y su ejecución numérica: deja a la vista, sobre el elemento real que el usuario construyó, los pasos intermedios que el software profesional encapsula. La literatura sobre enseñanza del método reporta que la simulación interactiva y la construcción paso a paso de las herramientas de cálculo favorecen la instrucción, con las salvedades que examina la \autoref{sec:fundamentos-pedagogicos}. Esa literatura orienta las decisiones de diseño de EduFEM.

En el plano práctico, EduFEM está concebido como apoyo a la cátedra. Es gratuito y de código abierto —su código fuente está disponible públicamente en línea\footnote{Repositorio del proyecto: \url{https://github.com/Sucullani/SoftwareED}}—, está íntegramente en español y organiza su interfaz según las tres fases habituales del análisis (pre-proceso, proceso y post-proceso), las mismas del software profesional. Además genera de forma automática una memoria de cálculo: un documento que reconstruye el procedimiento con sustitución numérica paso a paso, sobre los valores del modelo del propio usuario. Esta capacidad no se documenta en los antecedentes de software educativo revisados (\autoref{sec:antecedentes}). Con ella, el estudiante puede rehacer a mano, o en una hoja de cálculo, las operaciones que la memoria desarrolla —la matriz de deformación, la cuadratura de Gauss, la matriz de rigidez elemental y su ensamblaje— y contrastarlas con las del software, elemento por elemento.

En el plano técnico, construir un software propio en lugar de adoptar uno existente responde a la necesidad de controlar qué se expone del motor de cálculo. Un motor desarrollado sobre bibliotecas numéricas de código abierto \autocite{harris2020numpy,virtanen2020scipy} permite exhibir las matrices y las operaciones intermedias sin las restricciones de un producto cerrado, y verificar su exactitud frente a referencias reconocidas. La elección del lenguaje Python obedece a tres razones: no exige licencia al estudiante; reúne en un solo lenguaje el cálculo numérico, la interfaz gráfica y la generación de documentos; y permite distribuir el software como un instalador autónomo, que no requiere tener Python instalado.

En el plano profesional, el trabajo ofrece un ejemplo documentado de verificación y validación de un programa de cálculo estructural, el procedimiento del que depende la confianza en los resultados de cualquier programa de esa clase.

\section*{Objetivos}

El objetivo general de este trabajo es desarrollar EduFEM, un software educativo de elementos finitos que expone el procedimiento de cálculo de problemas bidimensionales de elasticidad lineal (tensión y deformación plana) con elementos cuadriláteros isoparamétricos Q4 y Q9, mediante su diseño e implementación en el lenguaje de programación Python y su verificación y validación frente a soluciones analíticas, casos clásicos de la literatura y un software comercial, para que el análisis por el Método de los Elementos Finitos resulte trazable y verificable paso a paso en la enseñanza de la Carrera de Ingeniería Civil de la Universidad Autónoma Tomás Frías.

Los objetivos específicos que concretan este propósito son los siguientes:

\begin{enumerate}
    \item Diagnosticar, a partir de los antecedentes y del estado del arte del software para la enseñanza del MEF, los atributos de trazabilidad y verificabilidad que los recursos disponibles no ofrecen, para fijar los requisitos de la herramienta.
    \item Fundamentar teóricamente el MEF para la elasticidad lineal plana con elementos isoparamétricos Q4 y Q9 —de la formulación variacional a la recuperación de tensiones, la calidad de malla, los fenómenos numéricos y los criterios de verificación y validación— y los principios que orientan la exposición del procedimiento, como base común del motor de cálculo, de los módulos educativos y de la memoria de cálculo.
    \item Desarrollar el motor de cálculo con elementos Q4 y Q9, la interfaz gráfica organizada en pre-proceso, proceso y post-proceso sobre un único lienzo, los módulos educativos que exponen cada etapa del procedimiento sobre el modelo real y la memoria de cálculo automática en PDF, con intercambio de datos (DXF, CSV), para que cada resultado pueda seguirse hasta sus datos.
    \item Verificar y validar el motor de cálculo mediante el método de soluciones manufacturadas y dos casos clásicos de la literatura (la viga de Timoshenko y la membrana de Cook), contrastando contra soluciones analíticas y un software comercial (SAP2000), para acotar el error de los resultados frente a referencias de exactitud conocida.
    \item Validar la propuesta contrastando la cobertura del procedimiento de cálculo y del contenido que la literatura especializada exige, junto con los resultados de la verificación, contra los criterios de aceptación fijados de antemano, para establecer si el análisis resulta trazable y verificable.
\end{enumerate}

\section*{Hipótesis y preguntas de investigación}

La \emph{hipótesis} de este trabajo es que el desarrollo de EduFEM permitirá mejorar la forma en que el software expone el procedimiento de cálculo —de caja negra a procedimiento a la vista—, elevando la trazabilidad y la verificabilidad del análisis por el Método de los Elementos Finitos en elasticidad plana en la Carrera de Ingeniería Civil de la Universidad Autónoma Tomás Frías.

En la nomenclatura de tres variables, la variable independiente es la forma en que el software expone el procedimiento de cálculo; la variable dependiente, la trazabilidad y la verificabilidad del análisis; y la variable interviniente, es decir, la solución tentativa con que se actúa sobre la primera, es EduFEM. Por tratarse de una investigación tecnológica, la hipótesis es la solución tentativa a un problema concreto, y se aprueba o no en la propia práctica de construir esa solución y ponerla a prueba \autocite[p.~85]{garciacordoba2005tecnologica}.

La hipótesis se comprueba con dos cláusulas, una por cada dimensión de la variable dependiente. Sus umbrales numéricos se fijan de antemano en la \autoref{sec:validacion-resultados}.

\begin{itemize}
    \item[(a)] \textbf{Trazabilidad.} Cada etapa del procedimiento de cálculo queda expuesta sobre el modelo del usuario, sea en un módulo educativo, en la memoria de cálculo o en el post-proceso. Las tres fases del análisis operan sobre un mismo lienzo. Y el contenido expuesto cubre las destrezas y los conceptos que un consenso publicado de expertos exige, dentro del alcance del trabajo.
    \item[(b)] \textbf{Verificabilidad.} Los resultados del motor coinciden, dentro de las tolerancias fijadas, con soluciones exactas construidas para la verificación, con la solución analítica de la viga de Timoshenko, con un modelo de SAP2000 y con el valor de referencia de la membrana de Cook. Reproducen además el bloqueo por cortante del elemento Q4 frente a la convergencia del Q9. Y la memoria de cálculo del ejemplo canónico puede reproducirse a mano.
\end{itemize}

Cada cláusula puede resultar falsa: la refutarían una etapa sin exposición, un contenido exigido sin instrumento, una tasa de convergencia fuera de margen, un error por encima del umbral o una memoria que no cierre con el cálculo manual. La hipótesis tampoco se cumple por definición. Exponer el procedimiento no asegura que lo expuesto esté completo ni que sea correcto, como muestra el defecto que este mismo trabajo detectó y corrigió en la extrapolación de tensiones a los nodos (\autoref{sec:leccion-extrapolacion}).

De esta hipótesis se desprenden dos preguntas de investigación, una por cláusula:

\begin{itemize}
    \item \textbf{PI-1} (trazabilidad). ¿Qué organización de la interfaz, de los módulos educativos y de la memoria de cálculo permite seguir cada resultado del análisis hasta sus datos, etapa por etapa, sobre el modelo del usuario, y cubre el contenido que la literatura especializada exige?
    \item \textbf{PI-2} (verificabilidad). ¿Reproduce el motor, dentro de tolerancias fijadas de antemano, las soluciones analíticas, los casos clásicos de la literatura y los resultados de un software comercial, así como los fenómenos numéricos que la teoría predice?
\end{itemize}

\section*{Metodología de la investigación}

Por su finalidad, este trabajo es \emph{propositivo}; por su naturaleza, es una \emph{investigación tecnológica}, la que busca conocimiento operativo y no solo establece cuál es la solución, sino que brinda la información para implementarla \autocite[p.~91]{garciacordoba2005tecnologica}. El método es el de la \emph{ciencia del diseño}, en la que el conocimiento sobre un problema y su solución se obtiene construyendo el producto tecnológico y evaluándolo con rigor \autocite[p.~75]{hevner2004design}. El trabajo sigue su proceso de seis actividades: identificar el problema, definir los objetivos de la solución, diseñar y desarrollar, demostrar, evaluar y comunicar \autocite[pp.~52-56]{peffers2007dsrm}. El enfoque es \emph{cuantitativo} y el nivel, \emph{descriptivo y comparativo}.

Este trabajo evalúa el software: su diseño, la exposición de su procedimiento de cálculo y la exactitud de sus resultados. Sus variables se miden sobre el software y sobre los casos de estudio. El diseño metodológico completo se desarrolla en la \autoref{sec:metodologia}.

\section*{Alcance y limitaciones}

El alcance de EduFEM está delimitado de manera explícita para lograr profundidad sobre un dominio acotado. En este documento, la expresión \emph{análisis estructural} se emplea en su acepción restringida de análisis de tensiones y deformaciones de elementos estructurales modelados como medios continuos bidimensionales, mediante la teoría de la elasticidad lineal. No comprende el análisis de sistemas estructurales de barras —pórticos o cerchas—, ni la formulación de placas, cáscaras o problemas tridimensionales.

La herramienta resuelve problemas \emph{bidimensionales} de \emph{elasticidad lineal estática}, en las dos idealizaciones clásicas de la mecánica de sólidos \autocite[p.~11]{timoshenko1970elasticity}: la \emph{tensión plana}, apropiada para cuerpos delgados cargados en su plano, y la \emph{deformación plana}, apropiada para cuerpos prismáticos largos cuya sección y cuya carga no varían a lo largo del eje. Dentro de ese dominio quedan problemas planos habituales en la ingeniería civil. En tensión plana, las vigas de gran canto y los muros cargados en su plano; en deformación plana, las secciones de presas, de muros de contención y de túneles, donde una rebanada representa todo el sólido.

La discretización se realiza con \emph{elementos cuadriláteros isoparamétricos} de cuatro nodos (Q4) y de nueve nodos (Q9), que cubren respectivamente la interpolación bilineal y la cuadrática completa y permiten contrastar el efecto del orden de interpolación sobre la precisión \autocites[pp.~137-138]{hughes2000fem}[pp.~164, 183-184]{onate2009structural}. La pareja Q4--Q9 es el par mínimo que contrasta dos órdenes de interpolación dentro de una sola familia de funciones de forma: el Q4 es bilineal y el Q9, bicuadrático. Ambos reproducen los polinomios lineales completos, de modo que la diferencia no está en la completitud lineal sino en el orden, que es lo que gobierna el comportamiento en flexión. Mantener una sola familia conserva además las matrices del seguimiento paso a paso en un tamaño legible. Los elementos triangulares y los cuadriláteros de ocho nodos quedan fuera de esta versión, y su incorporación se plantea en las recomendaciones. El producto es un \emph{software de escritorio} de carácter \emph{educativo}, desarrollado íntegramente en \emph{español}, con una interfaz diseñada bajo criterios de sencillez y claridad.

La elección del dominio bidimensional no es solo una restricción de alcance, sino una decisión de diseño. El continuo plano ocupa el punto intermedio entre el análisis unidimensional de barras y el análisis tridimensional de sólidos. El elemento de barra, con el que suele iniciarse la enseñanza del método, no alcanza para mostrar los conceptos centrales de la formulación isoparamétrica: el mapeo entre coordenadas naturales y físicas, la matriz Jacobiana, la matriz de deformación-desplazamiento y la cuadratura en dos direcciones solo aparecen al pasar al continuo. El análisis tridimensional sí contiene esos conceptos, pero el tamaño de sus matrices los oscurece: la matriz constitutiva crece a $6\times6$, la de deformación-desplazamiento a $6\times24$ en un hexaedro de ocho nodos y la rigidez elemental a $24\times24$ \autocite[p.~218]{cook2002concepts}. Esas dimensiones vuelven impracticable el seguimiento paso a paso que constituye el núcleo de la herramienta. El plano es la mínima dimensión en la que aparece el carácter de continuo y, a la vez, conserva matrices de tamaño legible —$\bm{D}$ de $3\times3$ y $\bm{B}$ de $3\times8$ en el elemento Q4—, lo que permite exhibir cada operación. El dominio plano actúa así como un puente: la formulación isoparamétrica, la cuadratura de Gauss y el ensamblaje que el software expone en dos dimensiones son los mismos que rigen el caso tridimensional \autocite{bathe2014fem,zienkiewicz2013fem}, planteado como línea de continuación del proyecto.

Las limitaciones fijadas de antemano son las siguientes. El software no aborda análisis tridimensional, comportamiento no lineal (del material o geométrico), efectos dinámicos ni dependientes del tiempo. El catálogo de materiales se restringe al comportamiento elástico, isótropo y lineal, definido por el módulo de elasticidad y el coeficiente de Poisson. La malla se construye manualmente —por tablas, por dibujo en el lienzo o por importación de un archivo DXF— o por generación estructurada en los casos de estudio: el software no incorpora mallado automático. La validación se realiza contra soluciones analíticas, casos clásicos de la literatura y un software comercial, no contra ensayos físicos. Estas fronteras responden a una decisión deliberada: privilegiar la transparencia y la solidez de un núcleo bien acotado por sobre la cobertura amplia de funcionalidades. Las limitaciones que el propio desarrollo puso de manifiesto se discuten en la \autoref{sec:interpretacion}.

\section*{Estructura del documento}

El documento se organiza en tres capítulos, además de esta introducción, de las conclusiones y de siete anexos.

El \autoref{cap:marco_teorico} desarrolla el marco teórico. Abre con los antecedentes y el estado del arte del software para la enseñanza del MEF, que sitúan el trabajo y cumplen el primer objetivo específico. Continúa con los principios que orientan la exposición del procedimiento. Desarrolla luego los fundamentos del método aplicados a la elasticidad lineal bidimensional, la calidad de malla y los criterios de verificación y validación, con lo que cumple el segundo objetivo.

El \autoref{cap:diseno} expone el diseño y la implementación. Abre con el diseño metodológico: el tipo de investigación y su método, las variables y su matriz de consistencia, la población y los casos de estudio, el procedimiento con sus instrumentos y los criterios de aceptación con que se contrasta la hipótesis. Continúa con el software, que materializa el tercer objetivo: los requisitos, las tecnologías empleadas, la arquitectura, el motor de cálculo, el pre-proceso, los módulos educativos, el post-proceso y la memoria de cálculo.

El \autoref{cap:resultados} presenta y analiza los resultados, con lo que cumple el cuarto y el quinto objetivos. Ejecuta la verificación y la validación numérica frente a soluciones analíticas y a un software comercial, mide la cobertura del procedimiento y del contenido, interpreta los hallazgos y contrasta la hipótesis.

Las conclusiones responden a cada objetivo en ese mismo orden, discuten el cumplimiento de la hipótesis y agrupan las recomendaciones por capítulo. Los anexos reúnen el manual de instalación (\autoref{anexo:instalacion}), el manual de uso (\autoref{anexo:manual}), una selección de listados de código del motor (\autoref{anexo:codigo}), los resultados extendidos de verificación y validación (\autoref{anexo:vyv}), el modelo de validación en SAP2000 con la solución analítica de Timoshenko (\autoref{anexo:sap2000}), los formatos de intercambio de datos (\autoref{anexo:formatos}) y una memoria de cálculo resuelta paso a paso sobre el ejemplo canónico (\autoref{anexo:memoria}).

\section*{Normas de citación y presentación}

Las citas y la lista de referencias siguen la norma Vancouver, en la versión del Comité Internacional de Editores de Revistas Médicas (ICMJE) y de la Biblioteca Nacional de Medicina de los Estados Unidos (NLM): numeración arábiga por orden de primera mención, títulos de revista completos y conjunción antes del último autor. El número va entre corchetes, y no en superíndice, para poder acompañarlo con la página citada. El localizador de página acompaña a toda cita que sostiene una ecuación, un valor numérico, un umbral, una definición o una atribución concreta, y se omite en las citas de encuadre. La presentación del documento sigue la séptima edición de las normas APA.
